\section{Annexe G -- Protocole V92}
Le protocole V92 transforme la confrontation V91 en une lecture d'article avec figures: on visualise $v_{\mathrm{info}}(z)$, les ratios de distance, et la séparation nette entre V90 et $\Lambda$CDM autour de la zone critique. Cette annexe est un module de présentation phénoménologique, pas une validation observationnelle indépendante.

\subsection{Objectif}
L'objectif de V92 est purement éditorial et analytique: produire un rendu directement intégrable dans l'article, avec deux graphes principaux et une interprétation claire des écarts relatifs au baseline standard. Les figures servent à lire un comportement modèle, pas à établir une déviation expérimentale finale.

\subsection{Figure 1: vitesse d'information}
La première figure compare la vitesse d'information V90 et la baseline $\Lambda$CDM sur la grille de test. Le point essentiel est que V90 est compressé autour de $z\approx0.46$ par le facteur géométrique local, alors que le baseline reste monotone et non corrigé.

\begin{figure}[htbp]
\centering
\includegraphics[width=0.92\linewidth]{v92_vinfo_comparison.png}
\caption{Comparaison de $v_{\mathrm{info}}(z)$ entre V90 et $\Lambda$CDM.}
\end{figure}

\subsection{Figure 2: ratios V90 / $\Lambda$CDM}
La seconde figure montre les ratios $d_{V90}/d_{\Lambda\mathrm{CDM}}$ et $v_{\mathrm{info},V90}/v_{\mathrm{info},\Lambda\mathrm{CDM}}$. Elle synthétise l'effet combiné du ralentissement du temps propre et de la compression géométrique.

\begin{figure}[htbp]
\centering
\includegraphics[width=0.92\linewidth]{v92_ratio_comparison.png}
\caption{Ratios de comparaison V90 / $\Lambda$CDM.}
\end{figure}

\subsection{Lecture scientifique}
Le message de V92 est simple: la zone critique ne se contente pas de modifier un paramètre global, elle produit un profil local sur $v_{\mathrm{info}}(z)$ et sur les ratios de distance. On voit ainsi trois effets distincts: le temps propre croît avec $z$, la vitesse d'information relative chute au voisinage de $z\approx0.46$, et les ratios V90 / $\Lambda$CDM restent inférieurs à 1 sur toute la grille testée. Cette lecture reste toutefois interne au modèle effectif et doit être confrontée à des données publiques pour devenir une affirmation physique.

\subsection{Résumé}
V92 est donc la version d'article de V91: même logique calculatoire, mais avec des graphes et un langage de lecture directement exploitable dans le manuscrit principal.